\subsection{Convolutions and Fundamental Solutions}\label{sec:convolutions-fundamental-solutions}

The classical definition of the convolution of two functions defined on $\mathbb{R}^m$ is
\begin{equation}\label{eq:classical-convolution-definition}
f*g(\bm{x})=\int_{\mathbb{R}^m}f(\bm{x}-\bm{y})g(\bm{y})\,\mathrm{d}\bm{y}.
\end{equation}
In this section, we shall consider a generalization of the definition to generalized functions and we shall give applications to the solution of partial differential equations with constant coefficients.

\subsubsection{The Direct Product of Distributions}\label{subsec:distribution-direct-product}

In general, one cannot define the product of two generalized functions $f(\bm{x})$ and $g(\bm{x})$. However, it is always possible to multiply two generalized functions depending on different variables. That is, if $f\in\mathcal{D}^{\prime}(\mathbb{R}^p)$ and $g\in\mathcal{D}^{\prime}(\mathbb{R}^q)$, then $f(\bm{x})g(\bm{y})$ can be defined as a distribution on $\mathbb{R}^{p+q}$.

\begin{Definition}
\label{def:distribution-direct-product}
Let $f\in\mathcal{D}^{\prime}(\mathbb{R}^p)$, $g\in\mathcal{D}^{\prime}(\mathbb{R}^q)$. Then the \textbf{direct product} $f(\bm{x})g(\bm{y})$ is the distribution on $\mathbb{R}^{p+q}$ given by
\begin{equation}\label{eq:distribution-direct-product-definition}
(f(\bm{x})g(\bm{y}),\phi(\bm{x},\bm{y}))=(f(\bm{x}),(g(\bm{y}),\phi(\bm{x},\bm{y}))).
\end{equation}

\end{Definition}

That is, we first regard $\phi(\bm{x},\bm{y})$ as a function only of $\bm{y}$, which depends on $\bm{x}$ as a parameter. To this function we apply the functional $g$. The result is then a real-valued function $\psi(\bm{x})$, which obviously has compact support. It is easy to show that $\psi$ is of class $C^\infty$ (Problem~\ref{ex:distribution-smooth-test-parameter-map}). Hence $\psi$ is a test function and $(f,\psi)$ is well defined. To justify the definition, it remains to be shown that $(f(\bm{x}),(g(\bm{y}),\phi_n(\bm{x},\bm{y})))$ converges to zero if $\phi_n$ converges to zero in $\mathcal{D}(\mathbb{R}^{p+q})$. Since $f$ is a distribution, it suffices to show that $\psi_n:=(g(\bm{y}),\phi_n(\bm{x},\bm{y}))$ converges to zero in $\mathcal{D}(\mathbb{R}^p)$. If $S_p\times S_q$ is a compact set containing the supports of all the $\phi_n$, then $S_p$ contains the supports of all the $\psi_n$. It remains to be shown that $\psi_n$ and all its derivatives converge uniformly to zero. Let $\alpha$ be a $p$-dimensional multi-index and let $\beta=(\alpha,0,\ldots,0)$. Assume that $D^\alpha\psi_n$ does not converge uniformly to zero. After choosing a subsequence, we may assume that there is a sequence of points $\bm{x}_n$ such that
\begin{equation}\label{eq:direct-product-continuity-contradiction}
|D^\alpha\psi_n(\bm{x}_n)|=|(g(\bm{y}),D^\beta\phi_n(\bm{x}_n,\bm{y}))|\geq\epsilon.
\end{equation}
But since the $\phi_n$ converge to zero with all their derivatives, the same is true for the sequence $\chi_n(\bm{y}):=D^\beta\phi_n(\bm{x}_n,\bm{y})$. Hence $\chi_n$ converges to zero in $\mathcal{D}(\mathbb{R}^q)$ and therefore $(g,\chi_n)$ converges to zero, a contradiction.

\begin{Example}
\label{ex:direct-product-delta}
As a simple example of a direct product, we note that \[\delta(\bm{x})\delta(\bm{y})=\delta(\bm{x},\bm{y}).\]
\end{Example}

If $\phi(\bm{x},\bm{y})$ has the special form $\phi_1(\bm{x})\phi_2(\bm{y})$, we obtain
\begin{equation}\label{eq:direct-product-separable-test-pairing}
(f(\bm{x})g(\bm{y}),\phi(\bm{x},\bm{y}))=(f,\phi_1)(g,\phi_2).
\end{equation}


Linear combinations of the form $\phi_1(\bm{x})\phi_2(\bm{y})$ are actually dense in $\mathcal{D}(\mathbb{R}^{p+q})$. To see this, let $\phi$ have support in the set $Q:=\{|\bm{x}|\leq a,\ |\bm{y}|\leq a\}$. By the Weierstrass approximation theorem (see Section~\ref{subsec:weierstrass-approximation}), there is a sequence of polynomials which converges to $\phi$ uniformly on the set $Q^{\prime}:=\{|\bm{x}|\leq2a,\ |\bm{y}|\leq2a\}$. Moreover, the argument used in the proof of the theorem also shows that the derivatives of the polynomials converge uniformly to those of $\phi$ on $Q^{\prime}$. We can thus choose polynomials $p_n$ in such a way that on $Q^{\prime}$ we have
\begin{equation}\label{eq:direct-product-polynomial-derivative-approximation}
|D^\alpha p_n-D^\alpha\phi|\leq\dfrac1n,\quad\forall\,|\alpha|\leq n.
\end{equation}
Let now $b_1(\bm{x})$, $b_2(\bm{y})$ be fixed test functions which are equal to 1 for $|\bm{x}|\leq a$ ($|\bm{y}|\leq a$) and equal to 0 for $|\bm{x}|\geq2a$ ($|\bm{y}|\geq2a$). Then the sequence
\begin{equation}\label{eq:direct-product-cutoff-polynomial-approximation}
\phi_n(\bm{x},\bm{y}):=b_1(\bm{x})b_2(\bm{x})p_n(\bm{x},\bm{y})
\end{equation}
converges to $\phi$ in $\mathcal{D}(\mathbb{R}^{p+q})$.

This fact and continuity can be used to show properties of the direct product by verifying them only on the restricted set of test functions of the form $\phi_1(\bm{x})\phi_2(\bm{y})$. One immediate conclusion is that in the definition we can evaluate $f$ and $g$ in the opposite order, i.e., we also have
\begin{equation}\label{eq:direct-product-reversed-pairing}
(f(\bm{x})g(\bm{y}),\phi(\bm{x},\bm{y}))=(g(\bm{y}),(f(\bm{x}),\phi(\bm{x},\bm{y})));
\end{equation}
we express this fact by the suggestive notation
\begin{equation}\label{eq:distribution-direct-product-commutativity}
f(\bm{x})g(\bm{y})=g(\bm{y})f(\bm{x}).
\end{equation}


Another obvious property is the associative law
\begin{equation}\label{eq:distribution-direct-product-associativity}
f(\bm{x})(g(\bm{y})h(\bm{z}))=(f(\bm{x})g(\bm{y}))h(\bm{z}).
\end{equation}

